\section*{PHÂN LỚP SENTIMENT ANALYSIS VỚI YOUTUBE API: \\ THU THẬP, CHUẨN HÓA VÀ GÁN NHÃN DỮ LIỆU BÌNH LUẬN PHIM \\TÓM TẮT}

\begin{itemize}
    \item[-] \textbf{Vấn đề nghiên cứu:}
    
    Đề tài nghiên cứu và xây dựng hệ thống phân tích cảm xúc (Sentiment Analysis) tự động từ bình luận YouTube sử dụng YouTube Data API v3 và thư viện TextBlob. Hệ thống thực hiện 6 bước xử lý: thu thập dữ liệu, chuẩn hóa văn bản, gán nhãn cảm xúc, thực nghiệm mô hình, xây dựng ứng dụng demo, và phân tích sâu kết quả. Dữ liệu được thu thập từ video "Avengers: Endgame" (2019) với 100,000 bình luận.
    
    Báo cáo gồm 5 chương:
        \begin{itemize}
        \item[•] Chương 1 - Giới thiệu: Tổng quan về sentiment analysis, mục tiêu và phạm vi đề tài
        \item[•] Chương 2 - Thu thập và Chuẩn hóa: YouTube API, quy trình làm sạch 8 bước
        \item[•] Chương 3 - Gán nhãn: Phương pháp TextBlob, 3 mức độ tin cậy
        \item[•] Chương 4 - Thực nghiệm và Ứng dụng: Phân tích mối quan hệ đặc trưng, xây dựng demo
        \item[•] Chương 5 - Phân tích sâu và Kết luận: Insights, đánh giá và hướng phát triển
        \end{itemize}
        
    \item[-] \textbf{Các hướng tiếp cận:}
        \begin{itemize}
        \item[•] Sử dụng YouTube Data API v3 với 4 API keys dự phòng
        \item[•] Áp dụng quy trình chuẩn hóa văn bản 8 bước (loại bỏ URL, mentions, hashtags, ký tự đặc biệt)
        \item[•] Gán nhãn cảm xúc với TextBlob (Positive/Negative/Neutral)
        \item[•] Phân tích tương quan giữa độ dài, likes và sentiment score
        \item[•] Trực quan hóa dữ liệu với 6 biểu đồ (bar, pie, scatter, box plot)
        \item[•] Xây dựng ứng dụng demo phân tích văn bản mới
        \end{itemize}
        
    \item[-] \textbf{Cách giải quyết vấn đề:}
    
    Sử dụng Python với các thư viện Pandas, NumPy, TextBlob, Matplotlib và Seaborn. Thu thập 100,000 bình luận từ YouTube API, sau đó làm sạch bằng regex và string processing. Gán nhãn cảm xúc dựa trên polarity score của TextBlob với 3 mức độ tin cậy (High/Medium/Low). Phân tích mối quan hệ giữa các đặc trưng bằng correlation analysis, distribution analysis và regression analysis. Xây dựng ứng dụng console cho phép người dùng nhập văn bản và nhận kết quả phân tích ngay lập tức.
    
    \item[-] \textbf{Một số kết quả đạt được:}
        \begin{itemize}
        \item[•] Thu thập thành công 100,000 bình luận, sau chuẩn hóa còn 92,521 bình luận hợp lệ
        \item[•] Phân bố sentiment: Neutral 59.5\%, Positive 28.6\%, Negative 11.9\%
        \item[•] Bình luận Positive có độ dài trung bình cao nhất (85.3 ký tự)
        \item[•] Bình luận Positive nhận được nhiều likes hơn (trung bình 12.8 likes)
        \item[•] Tương quan yếu giữa độ dài và sentiment score (r = 0.15)
        \item[•] Xây dựng thành công ứng dụng demo với giao diện console thân thiện
        \item[•] Tạo được 6 biểu đồ phân tích chi tiết về phân bố và mối quan hệ dữ liệu
        \item[•] Hệ thống xử lý và phân tích văn bản trong thời gian thực (< 1 giây)
        \end{itemize}
\end{itemize}

\newpage
